\documentclass[a4paper,12pt]{article}

% --- Use XeLaTeX font support ---

\usepackage{fontspec}
\setmainfont{TeX Gyre Termes} % nice Times-like font, installed by texlive-core
\setsansfont{TeX Gyre Heros} % optional sans-serif
\setmonofont{Fira Code}       % optional monospaced font
% --- Math packages ---
\usepackage{amsmath, amssymb, amsthm}
\usepackage{mathtools}

% --- Page layout ---
\usepackage[margin=1in]{geometry}

% --- Better lists ---
\usepackage{enumitem}
% --- Hyperlinks for references ---
\usepackage[colorlinks=true, linkcolor=blue, urlcolor=blue]{hyperref}

% Header: fancy layout like the rovided image
\usepackage{fancyhdr}
\setlength{\headheight}{26pt} % increase if fancyhdr warns
\pagestyle{fancy}
\fancyhf{} % clear header/footer
\rhead{\begin{tabular}{r} Mt.: 3052737\end{tabular}}
\chead{\begin{tabular}{c}Lösungen EA 2\end{tabular}}
\lhead{\begin{tabular}{l}61112 Lineare Algebra SS26 \end{tabular}}
\renewcommand{\headrulewidth}{0.8pt}
\fancyfoot[C]{- \thepage\ -}
% Ensure the plain page style (used by \maketitle) uses the same header

% --- Line spacing ---
\usepackage{setspace} % provides \doublespacing and \onehalfspacing

% --- Optional solutions environment ---
\newenvironment{solution}{
  \par\noindent\textbf{Solution:}\quad
}{\hfill$\blacksquare$\par\vspace{1em}}

% % --- Title info ---
% \title{Aufgabe \#1}
% \author{Your Name}
% \date{\today}
% 
\begin{document}
\raggedright
% Enable double (2.0) line spacing for the document body
\doublespacing
\noindent\textbf{Aufgabe 2.3} \\

Es gibt eine Basis zum Unterraum $U$, die sich aufgrund des Basisergänzungsatzes zu einer Basis des Vekroraums $V$ ergänzen lässt:
$$B = (u_1, \ldots , u_t, v_{t+1}, \ldots, v_{n}).$$
Zum Vektorraum $V$ gibt es einen Dualraum $V^*$. Zur Basis B gibts eine eindeutige Dualbasis $B^*$:
$$B^* = (u_1^*, \ldots , u_t^*, v_{t+1}^*, \ldots, v_{n}^*).$$
Aufgrund der Eigenschaften der Dualbasis gilt für die "hinteren" Linearformen $v_k^*$ (mit $k \ge t+1$):\\
\begin{itemize}
  \item $v_k^*(u_j) = 0$ für alle $j = 1, \ldots, t$ (sie verschwinden auf den Vektoren von $U$)
  \item $v_k^*(v_l) = 1$ falls $k=l$, und $0$ sonst.
\end{itemize}
Sind die gesuchten Linearformen $\alpha_i$ genau die hinteren Vektoren der Dualbasis? \\
Dazu setzen wir für $i = 1, \ldots, n-t$:
$$\alpha_i := v_{t+i}^*.$$

Zu zeigen ist:
$$
\boxed{U = \bigcap_{i=1}^{n-t} \ker(v_{t+i}^*)}
$$
\begin{itemize}
 \item Richtung: $U \subseteq  \bigcap_{i=1}^{n-t} \ker(v_{t+i}^*)$: \\
Sei $x$ ein Vektor aus dem Unterraum $U$, $x \in U$. Da $(u_1, \ldots, u_t)$ eine Basis von $U$ ist, lässt sich $x$ als Linearkombination der Basisvektoren schreiben:
$x = \sum_{j=1}^t \lambda_j u_j$.
Anwenden der "hinteren" Elemente der Dualbasis auf den Vektor:
$$v_k^*(x)=v_k^*(\sum_{j=1}^t \lambda_ju_j)=0 \ \ \text{für} \ \  k \geq t+1,$$
wegen $v_k^*(u_j) = 0$ für alle $j = 1, \ldots, t$ (sie verschwinden auf den Vektoren von $U$). Also liegt $x$ im Kern jeder dieser Linearformen, $x \in \bigcap_{i=1}^{n-t} \ker(v_{t+i}^*)$. \\

 \item Richtung: $ \bigcap_{i=1}^{n-t}\ker(v_{t+i}^*) \subseteq U$: \\
 Sei $x \in \bigcap_{i=1}^{n-t}\ker(v_{t+i}^*)$, also $v_{t+i}^*(x) = 0  \ \forall \ i$. \\
 Jedes $x$ lässt sich schreiben als
 $$x=\sum_{i=1}^t \lambda_iu_i+\sum_{i=t+1}^n\lambda_iv_i$$
 Da $x$ im Durchschnitt der Kerne liegt, wissen wir, dass $v_j^*(x)=0$ für alle $j=t+1, \ldots n$.
 Anwenden der $v_i^*$ auf die Basisdarstellung von x:
 $$v_j^*(x) =v_j^*(\sum_{i=1}^t \lambda_iu_i+\sum_{i=t+1}^n\lambda_iv_i) =0$$
 Für die erste Summe sind die Indices immmer verschieden, es gilt also $v_j^*(u_i) = 0$ für alle $i$. Bei der zweiten Summe bleibt entsprechend das $\lambda$ mit $i=j$ stehen, also
 $$v_j^*(x) = \lambda_j \ \text{für} \ j=t+1, \ldots, n.$$
 Es folgt zwingend
 $$\lambda_j = 0 \ \text{für} \ j=t+1, \ldots, n.$$
Also sind alle Koeffizienten der fortgesetzen Basisvektoren null. Übrig bleibt
$$x = \sum_{i=1}^t \lambda_i u_i$$
Das ist genau die Linearkombination der Vektoren des Unterraums, d.h. $x \in U$.
\end{itemize}
Damit sind beide Richtungen bewiesen und die Behauptung $U = \bigcap_{i=1}^{n-t} \ker(\alpha_i)$ gilt.
\end{document}
